
\section{Results}
\label{sec:results}

Table~\ref{tab:results-summary} gives a compact view of the principal
measurements.  Two experiments are reported separately.  The first is the
cross-representation sweep, in which all polynomial representations are
measured in one benchmark campaign with the same Goldilocks arithmetic
backend.  The second is the backend experiment, in which specialized
Goldilocks arithmetic and Montgomery arithmetic are compared for the same
evaluation algorithm.  Absolute timings are not compared across these two
campaigns.

\begin{table}[ht]
\centering
\caption{Summary of the main measurements at \(N=2^{20}\).  Cross-basis
times come from one common benchmark sweep.  Backend ratios are the median
of five independent same-run comparisons.}
\label{tab:results-summary}
\begin{tabular}{lll}
\toprule
Experiment & Quantity & Result\\
\midrule
Cross-basis &
Boolean-kernel &
1.0560 ms\\
Cross-basis &
parallel monomial &
1.4116 ms\\
Cross-basis &
kernel / parallel monomial &
1.337\(\times\)\\
Backend &
Montgomery / special, kernel &
1.2777\(\times\)\\
Backend &
Montgomery / special, monomial &
1.0591\(\times\)\\
\bottomrule
\end{tabular}
\end{table}

The first result is that the ranking of field multipliers and the ranking
of complete evaluators are not the same question.  In the isolated
multiplication benchmark, specialized Goldilocks multiplication and
Montgomery multiplication have nearly identical central estimates:
\(1.9107\) ns and \(1.9162\) ns, respectively.  Their measurement intervals
overlap.  No meaningful advantage can therefore be assigned to either
method from the isolated multiplication benchmark alone.

The kernel fold behaves differently.  The specialized operation
\[
t(a+b)+b
\]
takes \(3.0292\) ns in the microbenchmark, compared with \(3.7952\) ns for
the Montgomery implementation.  The ratio is \(1.2529\times\).  The
difference appears only after the multiplication has been embedded in the
compound fold operation, where the form of the Goldilocks modulus can be
used together with the surrounding additions.

\subsection{Field Backend and Complete Evaluators}
\label{sec:backend-results}

To test whether the microbenchmark difference survives in the complete
evaluator, the four full \(N=2^{20}\) configurations were repeated in
five independent Criterion invocations.  Table~\ref{tab:backend-stability}
reports the central estimate from each invocation.

\begin{table}[ht]
\centering
\caption{Repeated full-evaluator measurements at \(N=2^{20}\), in
microseconds.  M/S denotes Montgomery time divided by specialized
Goldilocks time within the same invocation.}
\label{tab:backend-stability}
\begin{tabular}{rrrrrr}
\toprule
Run & \(K_S\) & \(K_M\) & \(M_S\) & \(M_M\) & Kernel M/S\\
\midrule
1 & 883.61 & 1172.40 & 1453.50 & 1531.00 & 1.327\\
2 & 989.03 & 1227.70 & 1559.40 & 1654.20 & 1.241\\
3 & 1080.90 & 1159.90 & 1537.00 & 1901.10 & 1.073\\
4 & 1036.00 & 1323.70 & 2208.30 & 1666.00 & 1.278\\
5 & 952.15 & 1277.60 & 1607.80 & 1702.80 & 1.342\\
\midrule
Median &
989.03 &
1227.70 &
1559.40 &
1666.00 &
1.278\\
\bottomrule
\end{tabular}
\end{table}

The kernel result is stable in sign across all five invocations.  The
Montgomery implementation is slower in every run, with a median same-run
ratio of
\[
1.2777\times.
\]
The absolute specialized-kernel estimates vary from
\(883.61\,\mu\mathrm{s}\) to
\(1080.90\,\mu\mathrm{s}\), which is why the comparison is made
within each run rather than by combining absolute times collected at
different moments.

The monomial backend comparison is less stable.  The same-run ratios range
from
\[
0.7544\times
\quad\text{to}\quad
1.2369\times,
\]
and one invocation reverses the ordering.  The data therefore do not
support a claim that either arithmetic backend is systematically faster for
the complete blocked monomial evaluator.  This contrasts with the kernel
case and suggests that the specialized advantage is tied to the fold
structure rather than to scalar field multiplication alone.

\subsection{Cross-Representation Evaluation}
\label{sec:cross-representation}

We next compare the polynomial representations using one common Goldilocks
backend.  At \(N=2^{20}\), the optimized Boolean-kernel evaluator requires
\[
1.0560\ \mathrm{ms},
\]
while the parallel monomial implementation requires
\[
1.4116\ \mathrm{ms}.
\]
The ratio is
\[
\frac{1.4116}{1.0560}
=
1.337.
\]
Equivalently, the Boolean-kernel evaluator uses
\[
25.2\%
\]
less wall time for this isolated evaluation primitive.

The serial Horner method takes
\[
9.0402\ \mathrm{ms},
\]
which is
\[
8.561\times
\]
the kernel time.  Evaluation from the roots-of-unity Lagrange
representation is more expensive for an arbitrary off-domain point:
the parallel implementation takes
\[
17.1990\ \mathrm{ms},
\]
and the serial implementation
\[
21.3890\ \mathrm{ms}.
\]
These numbers characterize only arbitrary single-point evaluation from the
given representation; they do not imply that evaluation-form
representations are generally inefficient in workloads that operate
directly on the evaluation domain.

\subsection{Scaling with the Polynomial Size}
\label{sec:scaling}

Figure~\ref{fig:scaling} reports the complete cross-basis sweep from
\(N=2^{10}\) to \(N=2^{20}\).  The vertical axis is logarithmic in time.
The figure plays the same role as the global result plots in the reference
performance study: it shows the complete set of principal algorithms before
the individual measurements are interpreted.

\begin{figure}[ht]
\centering
\setlength{\unitlength}{1mm}
\begin{picture}(154,72)
\put(16,8){\vector(1,0){102}}
\put(16,8){\vector(0,1){58}}
\put(2,64){\makebox(14,4)[r]{\scriptsize time}}
\put(116,2){\makebox(12,4)[l]{\scriptsize \(N\)}}
\put(20,7){\line(0,1){2}}
\put(16,1){\makebox(8,4)[c]{\scriptsize \(2^{10}\)}}
\put(38,7){\line(0,1){2}}
\put(34,1){\makebox(8,4)[c]{\scriptsize \(2^{12}\)}}
\put(56,7){\line(0,1){2}}
\put(52,1){\makebox(8,4)[c]{\scriptsize \(2^{14}\)}}
\put(74,7){\line(0,1){2}}
\put(70,1){\makebox(8,4)[c]{\scriptsize \(2^{16}\)}}
\put(92,7){\line(0,1){2}}
\put(88,1){\makebox(8,4)[c]{\scriptsize \(2^{18}\)}}
\put(110,7){\line(0,1){2}}
\put(106,1){\makebox(8,4)[c]{\scriptsize \(2^{20}\)}}
\put(15,15.1){\line(1,0){2}}
\put(1,13.6){\makebox(13,3)[r]{\scriptsize 10 us}}
\put(15,29.2){\line(1,0){2}}
\put(1,27.7){\makebox(13,3)[r]{\scriptsize 100 us}}
\put(15,43.3){\line(1,0){2}}
\put(1,41.8){\makebox(13,3)[r]{\scriptsize 1 ms}}
\put(15,57.3){\line(1,0){2}}
\put(1,55.8){\makebox(13,3)[r]{\scriptsize 10 ms}}
\put(20,10.0){\circle*{1.8}}
\put(38,18.3){\circle*{1.8}}
\put(56,26.0){\circle*{1.8}}
\put(74,28.6){\circle*{1.8}}
\put(92,35.1){\circle*{1.8}}
\put(110,43.6){\circle*{1.8}}
\qbezier(20,10.0)(29.0,14.15)(38,18.3)
\qbezier(38,18.3)(47.0,22.15)(56,26.0)
\qbezier(56,26.0)(65.0,27.3)(74,28.6)
\qbezier(74,28.6)(83.0,31.85)(92,35.1)
\qbezier(92,35.1)(101.0,39.35)(110,43.6)
\put(20,14.2){\circle{2.2}}
\put(38,22.7){\circle{2.2}}
\put(56,26.9){\circle{2.2}}
\put(74,31.2){\circle{2.2}}
\put(92,37.4){\circle{2.2}}
\put(110,45.4){\circle{2.2}}
\qbezier(20,14.2)(29.0,18.45)(38,22.7)
\qbezier(38,22.7)(47.0,24.799999999999997)(56,26.9)
\qbezier(56,26.9)(65.0,29.049999999999997)(74,31.2)
\qbezier(74,31.2)(83.0,34.3)(92,37.4)
\qbezier(92,37.4)(101.0,41.4)(110,45.4)
\put(18.9,13.200000000000001){\framebox(2.2,2.2){}}
\put(36.9,21.7){\framebox(2.2,2.2){}}
\put(54.9,30.4){\framebox(2.2,2.2){}}
\put(72.9,38.6){\framebox(2.2,2.2){}}
\put(90.9,47.1){\framebox(2.2,2.2){}}
\put(108.9,55.6){\framebox(2.2,2.2){}}
\qbezier(20,14.3)(29.0,18.55)(38,22.8)
\qbezier(38,22.8)(47.0,27.15)(56,31.5)
\qbezier(56,31.5)(65.0,35.6)(74,39.7)
\qbezier(74,39.7)(83.0,43.95)(92,48.2)
\qbezier(92,48.2)(101.0,52.45)(110,56.7)
\put(20,29.8){\circle*{1.8}}
\put(38,34.8){\circle*{1.8}}
\put(56,38.2){\circle*{1.8}}
\put(74,44.6){\circle*{1.8}}
\put(92,52.5){\circle*{1.8}}
\put(110,60.7){\circle*{1.8}}
\qbezier(20,29.8)(29.0,32.3)(38,34.8)
\qbezier(38,34.8)(47.0,36.5)(56,38.2)
\qbezier(56,38.2)(65.0,41.400000000000006)(74,44.6)
\qbezier(74,44.6)(83.0,48.55)(92,52.5)
\qbezier(92,52.5)(101.0,56.6)(110,60.7)
\put(20,19.8){\circle{2.2}}
\put(38,28.1){\circle{2.2}}
\put(56,36.6){\circle{2.2}}
\put(74,45.1){\circle{2.2}}
\put(92,53.5){\circle{2.2}}
\put(110,62.0){\circle{2.2}}
\qbezier(20,19.8)(29.0,23.950000000000003)(38,28.1)
\qbezier(38,28.1)(47.0,32.35)(56,36.6)
\qbezier(56,36.6)(65.0,40.85)(74,45.1)
\qbezier(74,45.1)(83.0,49.3)(92,53.5)
\qbezier(92,53.5)(101.0,57.75)(110,62.0)
\put(122,60){\circle*{1.8}}
\put(126,58.5){\makebox(24,3)[l]{\scriptsize \(K\)}}
\put(122,53){\circle{2.2}}
\put(126,51.5){\makebox(24,3)[l]{\scriptsize \(M_p\)}}
\put(120.9,44.9){\framebox(2.2,2.2){}}
\put(126,44.5){\makebox(24,3)[l]{\scriptsize \(M_h\)}}
\put(122,39){\circle*{1.8}}
\put(126,37.5){\makebox(24,3)[l]{\scriptsize \(L_p\)}}
\put(122,32){\circle{2.2}}
\put(126,30.5){\makebox(24,3)[l]{\scriptsize \(L_s\)}}
\end{picture}
\caption{Single-point evaluation time as a function of \(N\) for the five
implementations in the common cross-basis benchmark.  \(K\): optimized
Boolean-kernel; \(M_p\): parallel blocked monomial; \(M_h\): serial Horner;
\(L_p\): parallel Lagrange; \(L_s\): serial Lagrange.  The vertical
placement is logarithmic in measured time.}
\label{fig:scaling}
\end{figure}

The small-size regime is affected strongly by parallel scheduling and fixed
overheads.  At \(N=2^{10}\), the kernel evaluator takes approximately half
the time of the parallel monomial evaluator, but the absolute difference is
only a few microseconds.  At \(N=2^{14}\), the two parallel methods become
closer, with a ratio of approximately
\[
1.169\times.
\]
For larger inputs the separation grows again.  At \(N=2^{18}\), the
parallel monomial evaluator is
\[
1.447\times
\]
the kernel time, and at \(N=2^{20}\) the ratio is
\[
1.337\times.
\]

The serial Horner curve grows linearly with a long dependency chain and
separates rapidly from the two parallel implementations.  The two Lagrange
curves have the largest constants because arbitrary off-domain evaluation
includes denominator construction, batch inversion, and a full linear
accumulation.

\subsection{Interpretation}
\label{sec:results-interpretation}

The measurements support two separate conclusions.  First, the
Boolean-kernel representation gives a lower single-point evaluation time
than the fair parallel monomial baseline in the common \(N=2^{20}\)
benchmark.  Second, this advantage cannot be explained by a generally
faster scalar field multiplication: isolated Goldilocks and Montgomery
multiplication are nearly tied, whereas the specialized Boolean-kernel fold
and the complete specialized kernel evaluator show a clear advantage.

The result is therefore an interaction between representation and
implementation.  The Boolean-kernel basis exposes a balanced reduction
tree, and the Goldilocks modulus permits the arithmetic inside each fold to
be specialized as one compound operation.  The benchmark does not establish
a corresponding whole-prover speedup.  Whether the gain survives in a
proof system depends on how long the polynomial can remain in the
Boolean-kernel representation and on the cost of any required basis
conversion.
